\section{Benchmark Design}
\label{sec:design}
The benchmark has two tracks, one per family of soft binding, that share their principles, their hosts and, for the combination experiment, their media. This section states what they share and then describes each track's protocol.

\subsection{Shared Principles}
\label{sec:principles}
\paragraph{Licence audit.} A provenance service can deploy only models whose code \emph{and} weights permit commercial use, and the two licences differ in several projects; the weights may also have been trained on data restricted to non-commercial use. We recorded all three for every method. Methods that fail the audit were measured as reference points and are marked \dag\ (watermarks) or \reftier\ (fingerprints) in every table; they are excluded from rankings and recommendations.

\paragraph{Public, pinned media.} All media are public and were fetched from their distribution points; every file is pinned by SHA-256 in a manifest. No private or personal media were used. All randomness (payloads, transformations, splits) is seeded, so every method sees byte-identical inputs.

\paragraph{False matches are controlled, not just reported.} A soft binding is useful only if it rarely points unmarked or unregistered content at someone else's manifest. Both tracks therefore fix the decision rule before measuring robustness. For watermarks, the rule is either the model's blind detector or an expected-payload test that requires $k$ of $n$ decoded bits to match the key the verifier expects, with $k$ chosen for a per-key false-match probability of at most $10^{-6}$ under independent fair bits. For fingerprints, thresholds are fixed on one half of the never-registered negatives, at query-level and at registry-size-independent pair-level false-match rates, and evaluated on the other half; a target is reported only if the calibration set resolves it. Detection is always reported at such an operating point, next to threshold-free measures.

\paragraph{The source is the unit.} Every image, clip or track contributes one query per transformation, and the transformed copies of one source succeed or fail together far more often than independent queries would. Intervals therefore come from bootstraps over sources (for fingerprint retrieval rates and paired differences, with the threshold recomputed from the resampled calibration set in every replicate), exact intervals for watermark false positives count sources rather than trials, and paired comparisons pair by source. For audio watermarks, whose clips were rendered at three sample rates, the clip is the source.

\paragraph{Hosts.} GPU work ran on NVIDIA A10G instances (AWS g5, AMD EPYC 7R32) and CPU work on the same processor family; latency is reported per network on an idle host. Because the watermarking track found host-dependent decoder outputs, both tracks were also run or decoded on an Apple M5 Pro (arm64) laptop to test determinism.

\subsection{Watermarking Protocol}
\label{sec:wm-method}
Fig.~\ref{fig:wm-pipeline} summarises the protocol. Each configuration embeds a pseudo-random payload into every corpus item, the marked item is measured against the original, a fixed set of transformations is applied, and the decoder output is compared with the embedded payload. Unmarked items pass through the same decoders to measure false positives. All randomness is seeded, and payloads are derived from a hash of a fixed seed and the item index, so every configuration of a given payload size receives the same message for the same item.

\begin{figure*}[!t]
\centering
% Benchmark protocol schematic (Fig. 1). Drawn in TikZ so text matches the document.
\begin{tikzpicture}[
  font=\sffamily\footnotesize,
  node distance=4.5mm and 5mm,
  box/.style={draw=black!70, line width=0.5pt, rounded corners=1.5pt, align=center, minimum height=9mm,
              inner xsep=2.5pt, fill=white},
  data/.style={box, fill=black!5},
  metric/.style={box, draw=none, fill=blue!8, text width=26mm, minimum height=8mm},
  arr/.style={-{Stealth[length=1.8mm]}, line width=0.5pt, draw=black!70},
  lab/.style={font=\sffamily\scriptsize, text=black!65, align=center},
]
% Top row: marked path
\node[data] (src) {Public corpus\\\scriptsize SHA-256 pinned};
\node[box, right=of src] (emb) {Embed\\\scriptsize seeded payload};
\node[data, right=of emb] (marked) {Marked item};
\node[box, right=of marked] (att) {Transform\\\scriptsize \val{wm:count/transforms} settings};
\node[box, right=of att] (dec) {Decode};
\node[box, right=of dec] (cmp) {Compare with\\\scriptsize expected payload};

\draw[arr] (src) -- (emb);
\draw[arr] (emb) -- (marked);
\draw[arr] (marked) -- (att);
\draw[arr] (att) -- (dec);
\draw[arr] (dec) -- (cmp);

% Bottom row: unmarked path for false positives
\node[box, below=9mm of att] (att0) {Benign transform\\\scriptsize or none};
\node[box, right=of att0] (dec0) {Decode};
\node[box, right=of dec0] (fp) {Blind detection;\\\scriptsize match vs.\ $k$ of $n$};
\draw[arr] (src.south) |- (att0.west);
\draw[arr] (att0) -- (dec0);
\draw[arr] (dec0) -- (fp);
\node[lab, anchor=south west] at ($(src.south |- att0.west)+(1.5mm,0.6mm)$) {unmarked items};

% Measurements
\node[metric, above=6mm of marked] (mq) {\textbf{Quality}\\\scriptsize FLIP, LPIPS, PESQ, VMAF};
\node[metric, above=6mm of cmp] (mr) {\textbf{Robustness}\\\scriptsize bit accuracy, worst subset};
\node[metric, below=6mm of fp] (mf) {\textbf{False positives}\\\scriptsize exact 95\% intervals};
\node[metric, above=6mm of emb, xshift=-16mm] (mt) {\textbf{Latency, cost}\\\scriptsize GPU and CPU host};
\draw[arr, densely dashed] (marked) -- (mq);
\draw[arr, densely dashed] (cmp) -- (mr);
\draw[arr, densely dashed] (fp) -- (mf);
\draw[arr, densely dashed] (emb.north) -- ++(0,2.5mm) -| (mt.south);

% Qualification
\node[box, fill=black!85, text=white, draw=none, left=10mm of mf, text width=44mm] (q)
  {\textbf{Qualification}\\\scriptsize open licence, gate bit accuracy $\ge 0.95$\\\scriptsize (pooled and worst subset), false-positive evidence};
\draw[arr] (mf) -- (q);
\end{tikzpicture}

\caption{Benchmark protocol. Each of the \val{wm:count/image} image, \val{wm:count/audio} audio and \val{wm:count/video} video configurations embeds a seeded payload into every item. Quality is measured on the marked item; robustness on the decoder output after each transformation; false positives on unmarked items under the untouched condition and three benign transformations; latency and cost on a fixed GPU and CPU host. A configuration qualifies for ranking only if its licence, gate robustness and false-positive evidence all pass.}
\label{fig:wm-pipeline}
\end{figure*}

\subsubsection{Methods Under Test}
We selected watermarking methods that are distributed with trained weights and an inference implementation, and that target natural images, speech and music, or natural video. Table~\ref{tab:wm-models} lists them. For each method we recorded the licence of the code and, separately, of the weights, because the two differ in several projects. A configuration counts as \emph{open} only when both carry a licence that permits commercial use. Methods whose weights have no stated licence, or whose weight provenance we could not establish, were still measured and are reported as reference points, but they are excluded from every ranking. Several methods were excluded before measurement: VINE, whose code licence is non-commercial \cite{lu2025vine}; the WAM checkpoint trained on COCO, licensed CC BY-NC \cite{sander2025wam}; and DVMark, for which no weights are released \cite{luo2025dvmark}.

Strength was varied where a method exposes a single scalar that scales the embedded residual: the TrustMark strength factor, the blending factor of the PixelSeal, Video Seal and WAM embedders, and the AudioSeal output scale. We kept each method's default and added settings on both sides of it. For PixelSeal we added two settings between 0.1 and the default 0.2 after an initial pass showed that the robustness threshold is crossed in that interval. All other hyperparameters were left at the values published with the weights.

\begin{table}[!t]
\centering
\footnotesize
\caption{Methods, payload sizes and licences.}
\label{tab:wm-models}
\setlength{\tabcolsep}{3pt}
\begin{tabular}{llrll}
\toprule
Modality & Method & Bits & Code & Weights \\
\midrule
Image & TrustMark B, C, P, Q \cite{bui2025trustmark} & 100 & MIT & MIT \\
Image & PixelSeal \cite{soucek2025pixelseal} & 256 & MIT & MIT \\
Image, video & Video Seal 1.0 \cite{fernandez2024videoseal} & 256 & MIT & MIT \\
Image, video & ChunkySeal \cite{petrov2025chunkyseal} & 1024 & MIT & MIT \\
Image & WAM \cite{sander2025wam} & 32 & MIT & MIT$^{a}$ \\
Image & DWT-DCT, DWT-DCT-SVD \cite{invisiblewatermark} & 64 & MIT & n/a \\
Image & RivaGAN \cite{zhang2019rivagan} & 32 & MIT & unclear\textsuperscript{\dag} \\
Image & InvisMark \cite{xu2025invismark} & 100 & MIT & none stated\textsuperscript{\dag} \\
Audio & AudioSeal base, streaming \cite{sanroman2024audioseal} & 16 & MIT & MIT \\
Audio & WavMark \cite{chen2023wavmark} & 16 & MIT & MIT \\
Audio & SilentCipher \cite{singh2024silentcipher} & 40 & MIT & none stated\textsuperscript{\dag} \\
Audio & Perth \cite{perth2025} & 0 & MIT & MIT \\
\bottomrule
\end{tabular}
\par\vspace{2pt}\parbox{\linewidth}{\raggedright\scriptsize Bits: payload bits carried and compared. TrustMark transmits a 100-bit codeword; in the BCH\_5 mode used here it carries a 61-bit payload, 35 BCH parity bits and 4 version bits. Bit accuracy and the expected-payload test use the 100 transmitted bits; exact recovery compares the 61 decoded payload bits. $^{a}$The MIT-licensed checkpoint trained on SA-1B; the checkpoint trained on COCO is licensed CC BY-NC and was not used. \dag~Reported for reference only. Perth signals presence without a payload. Every checkpoint is pinned by SHA-256 in the lockfile kept with the scripts (Data and Code Availability).}
\end{table}

\subsubsection{Corpus}
\emph{Images} (\val{wm:image/n_items} items): the 24 Kodak images \cite{kodak}, 25 DIV2K validation images \cite{agustsson2017div2k} at about 2K resolution, 28 CLIC 2020 professional validation images \cite{clic}, and 8 CC0 high-dynamic-range panoramas from Poly Haven \cite{polyhaven}, exposure-normalised to a median luminance of 0.18 and stored as 16-bit sRGB TIFF. DIV2K images are used for measurement only and are not redistributed. FLIP was computed on the 60 Kodak, CLIC and HDR images and ColorVideoVDP on the 32 Kodak and HDR images, because both are costly at DIV2K resolution; the other metrics use all 85 images.

\emph{Audio} (\val{wm:audio/n_items} inputs from \val{wm:audio/n_sources} source clips): 40 LibriSpeech test-clean utterances from 40 different speakers \cite{panayotov2015librispeech} and 7 public-domain or CC BY music excerpts distributed with librosa \cite{mcfee2015librosa}, each resampled to 16, 44.1 and 48~kHz. One music excerpt failed in every configuration at all three rates; we treat a clip that fails for every method as a corpus defect and exclude it everywhere, leaving \val{wm:audio/n_sources} source clips (40 speech, 6 music) and 138 inputs, of which \val{wm:audio/n_speech} are speech and \val{wm:audio/n_music} music. The three renderings of a clip are not independent samples; Section~\ref{sec:wm-stats} explains how the statistics account for this.

\emph{Video} (\val{wm:video/n_items} clips): the first 2.5~s of six 1080p sequences from the Xiph test-media collection \cite{xiphderf} (\emph{aspen}, \emph{controlled burn}, \emph{red kayak}, \emph{speed bag}, \emph{crowd run} and the \emph{Sintel} trailer), covering slow and fast motion, fine texture and synthetic content.

\subsubsection{Transformations}
Image transformations (\val{wm:count/transforms/image}): JPEG at quality 90, 75, 60 and 50; WebP at 80 and 60; bicubic downscaling to 0.75 and 0.5; centre crops retaining 90, 70 and 50\% of each side; Gaussian blur with $\sigma=1$; and brightness scaling by 0.8 and 1.2. Audio transformations (\val{wm:count/transforms/audio}): MP3 at 320, 128 and 64~kb/s; AAC at 256, 128 and 64~kb/s; resampling round trips through 22.05 and 8~kHz; gain of $-6$~dB; and additive white noise at 30~dB SNR. Codec round trips are aligned to the original by cross-correlation before decoding, so that encoder delay does not register as watermark loss. A verifier that holds only the received file cannot do this, so a follow-up experiment decoded the same round trips without alignment (Section~\ref{sec:wm-res-audio}). Video transformations (\val{wm:count/transforms/video}): H.264 at CRF 18, 22, 23 and 28; H.265 at CRF 23 and 28; and H.264 at CRF 23 after 0.5 downscaling or a 75\% centre crop. Marked video is stored losslessly (FFV1, 4:4:4) before each transformation, so the transformation is the only lossy step.

One transformation per modality serves as the \emph{gate}: JPEG 75 for images, MP3 at 128~kb/s for audio and H.264 at CRF 23 for video. These are common default settings in sharing pipelines, and a watermark that does not survive them is of little use for provenance.

\subsubsection{Perceptual Quality}
Image quality was measured with PSNR, SSIM \cite{wang2004ssim}, MS-SSIM \cite{wang2003msssim}, LPIPS \cite{zhang2018lpips}, FLIP \cite{andersson2020flip} and ColorVideoVDP \cite{mantiuk2024cvvdp}. We use FLIP as the primary image metric because it models the visibility of differences when flipping between two images, which is the viewing condition under which an embedded pattern is most likely to be noticed. Audio quality was measured with SNR, SI-SNR \cite{leroux2019sdr}, wideband PESQ \cite{itu2007p8622}, STOI \cite{taal2011stoi}, the absolute change in integrated loudness following ITU-R BS.1770 \cite{itu2023bs1770}, and a log-spectral difference. PESQ is primary. It was designed and validated for speech, so the primary PESQ value is the mean over the \val{wm:audio/n_speech} speech inputs only; PESQ on the \val{wm:audio/n_music} music inputs is reported separately and not used for ranking. Video quality was measured with PSNR, SSIM, LPIPS and FLIP averaged over frames, and with VMAF \cite{rassool2017vmaf}, which is primary.

\subsubsection{Robustness}
Robustness is the fraction of payload bits recovered correctly after a transformation (bit accuracy), averaged over items. Because a verifier acts on a decision rather than on a bit rate, we also report two payload-level rates (Table~\ref{tab:wm-exact}): \emph{exact recovery}, the share of items whose decoded payload equals the embedded one, after BCH decoding of its 61-bit payload for TrustMark (the only method with an error-correcting layer) and with every bit correct for the others; and the share of items that pass the expected-payload test of Section~\ref{sec:wm-fp-method} ($\geq k$ of $n$ bits). For video, decoder logits are averaged over all frames of a clip before thresholding, which is the aggregation the Video Seal and PixelSeal decoders are designed for. We report the pooled mean and the mean on the worst content subset (image corpus, audio source and sample rate, or video clip), because a pooled mean can hide systematic failure on one kind of content.

\subsubsection{False Positives}
\label{sec:wm-fp-method}
We distinguish two decision rules. A \emph{blind} detector uses only the model's own output: TrustMark's error-correcting decoder reports whether a valid codeword was found, and the audio models expose a detection score, which we threshold at 0.8 for AudioSeal (whose score is the fraction of audio frames classified as watermarked) and at 0.5 for the others. An \emph{expected-payload} test instead declares a watermark when at least $k$ of $n$ decoded bits equal the payload the verifier expects. Under the null hypothesis that unmarked content yields independent fair bits, the per-key false-match probability is
\begin{equation}
P_{\mathrm{FM}}(k,n) = \sum_{i=k}^{n}\binom{n}{i}2^{-n},
\end{equation}
and we set $k$ to the smallest value with $P_{\mathrm{FM}}\le 10^{-6}$. For 16-bit payloads no such $k$ exists, since even an exact match has probability $2^{-16}\approx1.5\times10^{-5}$; for these methods only the blind detector can be tested.

False-positive trials run each decoder on the unmarked corpus under no transformation and three benign ones (JPEG 75, 0.5 downscale and 70\% crop for images; the untouched signal, MP3 at 128~kb/s and a 22.05~kHz round trip for audio; and, for video, the H.264 CRF~18 encode of each unmarked clip). A verifier checks one fixed expected payload, the one recorded for the asset; comparing each decode with the item's own payload is that check. We also compare each decode with 50 unrelated payloads (200 for video), one fixed list shared by all items. These comparisons test whether decoded bits agree with payloads the content was never marked with, conditional on the decodes we have; they reuse the same decodes and therefore add no independent media, and we report them as counts and as the largest number of matching bits observed, not as additional trials. The transformed copies of one item, and the three renderings of one audio clip, share their source, so observed rates and their exact Clopper-Pearson intervals \cite{clopper1934} use the unmarked source as the trial unit: \val{wm:fpr/image/PixelSeal/own_src_n} images, \val{wm:fpr/audio/WavMark/blind_src_n} audio clips and \val{wm:fpr/video/PixelSeal/own_src_n} video clips. All intervals are two-sided at 95\%. We also plot the distribution of matching bits against the binomial null.

A configuration's false-positive evidence is \emph{analytic} when an expected-payload test with $P_{\mathrm{FM}}\le10^{-6}$ exists and no unmarked trial reached $k$, and \emph{observed} when no such test exists but the blind detector produced no false detection. Our target is a rate of at most $10^{-3}$. Table~\ref{tab:wm-fpr} keeps the two kinds of evidence apart: the theoretical $P_{\mathrm{FM}}$ of each expected-payload test, and the observed counts with their source-level 95\% upper bounds. Analytic evidence meets the target by construction, provided unmarked content decodes to bits that are independent of the expected payload; the match distributions in Section~\ref{sec:wm-res-fpr} are consistent with that assumption but cannot prove it. Observation alone does not establish the target for any configuration: with no false positive, the upper bound falls below $10^{-3}$ only after \val{wm:fpr/sources_for_1e-3} independent sources, far more than our corpus holds.

\subsubsection{Latency and Cost}
\label{sec:wm-cost}
Embedding latency was measured on an AWS g5.xlarge instance (NVIDIA A10G, 24~GB; AMD EPYC 7R32, 4 vCPU; PyTorch 2.14, CUDA 13.0) as the median over all corpus items, excluding the first item of each run (which includes lazy initialisation), with device synchronisation before and after each call. We report milliseconds per megapixel for images and video frames, and milliseconds of compute per second of audio. CPU latency was measured on the same host with four threads, as the median over four images (Kodak and DIV2K) or audio items after one warm-up item. Model loading time is excluded.

Strength only scales the residual that the network adds, so settings of one network perform identical computation. Repeated settings nevertheless differed by up to about 20\% in measured latency, which is timing noise. We therefore assign each configuration the median latency of its network, and treat differences between networks of less than about 10\% as unresolved.

Cost is an estimate of on-demand compute price per unit of media:
\begin{equation}
C = \frac{t \cdot U}{3.6\times10^{6}}\, p,
\end{equation}
where $t$ is the latency in milliseconds per unit, $U$ the number of units (1{,}000 megapixels, or 1{,}000 hours of audio expressed in seconds), and $p$ the hourly price: \$\val{wm:price/gpu} for g5.xlarge and \$\val{wm:price/cpu} for m6a.xlarge, a 4-vCPU instance of the same processor family, both us-east-1 on-demand Linux prices effective \val{wm:price/date} and retrieved from the AWS Price List API. Each configuration is costed on the cheaper of the two. The estimate assumes a fully utilised instance and excludes decoding, storage, data transfer and idle capacity; it is intended for comparing methods, not for budgeting a deployment.

\subsubsection{Qualification and Combined Score}
A configuration \emph{qualifies} when (i) its code and weights are openly licensed, (ii) gate bit accuracy is at least 0.95 both pooled and on the worst subset, and (iii) its false-positive evidence is analytic or observed. The best-cost, best-speed and best-combined configurations in each table are chosen among qualifying configurations only: a fast method whose mark does not survive the gate transformation is not a usable operating point.

The combined score summarises four criteria as ratios to the best qualifying configuration, so that each lies in $(0,1]$ and equals 1 for the best:
\begin{align}
r_q &= \begin{cases} F^{*}/F & \text{image (FLIP } F\text{)}\\ (P-1)/(P^{*}-1) & \text{audio (PESQ } P\text{)} \\ V/V^{*} & \text{video (VMAF } V\text{)}\end{cases}\nonumber\\
r_b &= (a-0.5)/(a^{*}-0.5),\quad r_t = t^{*}/t,\quad r_c = C^{*}/C,\nonumber\\
S &= \exp\Big(\sum_{j} w_j \ln r_j\Big),\quad \sum_j w_j = 1,
\end{align}
where $a$ is worst-subset gate bit accuracy (0.5 is chance) and starred quantities are the best qualifying values. With equal weights, $S$ is the geometric mean of the four ratios, which rewards balance and penalises a configuration that is poor on any one criterion. Because speed and cost are strongly correlated, we also report a quality-weighted variant ($w_q=0.5$, $w_b=0.2$, $w_t=w_c=0.15$) and note where the two disagree.

\subsubsection{Statistics}
\label{sec:wm-stats}
The experimental unit is the source: an image, a video clip, or an audio clip together with its three sample-rate renderings. Means are over items, and their percentile bootstrap 95\% intervals come from 10{,}000 resamples of sources, so that the three renderings of an audio clip are drawn together and the intervals reflect \val{wm:audio/n_sources} independent audio sources rather than 138 inputs (Table~\ref{tab:wm-ci} in the Appendix). Planned comparisons between the leading configurations use two-sided Wilcoxon signed-rank tests on per-source differences \cite{wilcoxon1945}, averaging the renderings of an audio clip first, with Holm adjustment over the family of five tests \cite{holm1979}. Video results rest on six clips and are therefore descriptive; we do not claim significance for video differences.

\subsubsection{Cross-Device Determinism}
To test whether decoding is reproducible across hardware, we decoded the same compressed files on the benchmark host (x86-64, Linux) and on an Apple M-series machine (arm64, macOS), each with one and four CPU threads and with its accelerator (CUDA or MPS). Files were twelve Video Seal-marked clips (three clips, strengths 0.4 and 1.1, CRF 18 and 22) and one TrustMark-Q-marked image after JPEG 75. For each file we recorded the largest logit difference within one host and across hosts, the number of flipped bits, and whether decoded frames were byte-identical across hosts before and after conversion from YUV to RGB. To separate the colour conversion from the network, each host then converted the first \val{wm:rgb/nframes} frames of every clip to RGB once and saved the arrays, and both hosts decoded both hosts' arrays with the same TorchScript detector (a 2$\times$2 design: input host by decoding host). Both hosts used PyAV \val{wm:rgb/av} with libswscale \val{wm:rgb/swscale}.


\subsection{Fingerprinting Protocol}
\label{sec:fp-method}
Fig.~\ref{fig:fp-protocol} summarises the protocol. Each method fingerprints a set of registered assets and a larger set of distractors, which together form the registry. Registered assets are then transformed or partially edited and used as queries; media that were never registered go through the same transformations and serve as negative queries. Every query is searched exhaustively against the registry, and a threshold fixed on one half of the negatives decides whether the best match is reported. All randomness is seeded, so every method sees byte-identical queries.

\begin{figure*}[!t]
\centering
% Benchmark protocol schematic (Fig. 1). Drawn in TikZ so text matches the document.
\begin{tikzpicture}[
  font=\sffamily\footnotesize,
  box/.style={draw=black!70, line width=0.5pt, rounded corners=1.5pt, align=center, minimum height=9mm,
              text width=30mm, inner xsep=2pt, fill=white},
  data/.style={box, fill=black!5},
  metric/.style={box, draw=none, fill=blue!8, minimum height=10mm, text width=33mm},
  arr/.style={-{Stealth[length=1.8mm]}, line width=0.5pt, draw=black!70},
  lab/.style={font=\sffamily\scriptsize, text=black!60},
]
% lane 1: registration
\node[data] (reg) at (0,0) {Registered assets\\\scriptsize images, audio, video};
\node[box] (wm) at (3.85,0) {Optional watermark\\\scriptsize TrustMark, AudioSeal,\\[-2pt]\scriptsize Video Seal};
\node[box] (fp) at (7.7,0) {Fingerprint\\\scriptsize \val{fp:count/methods} methods};
\node[data] (idx) at (11.55,0) {Registry index\\\scriptsize assets + distractors};
\draw[arr] (reg) -- (wm);
\draw[arr] (wm) -- (fp);
\draw[arr] (fp) -- (idx);
\node[lab, anchor=south] at (5.775,0.72) {registration};

% lane 2: queries
\node[data] (neg) at (0,-1.9) {Never-registered\\\scriptsize media (negatives)};
\node[box] (att) at (3.85,-1.9) {Transform or edit\\\scriptsize \val{fp:count/attacks} attacks, 6 edits};
\node[box] (qfp) at (7.7,-1.9) {Fingerprint\\\scriptsize query};
\node[box] (srch) at (11.55,-1.9) {Exact search\\\scriptsize top-50, threshold};
\draw[arr] (reg.south east) -- node[lab, right, pos=0.45, inner sep=3pt] {positives} (att.north west);
\draw[arr] (neg) -- node[lab, below, inner sep=2pt] {negatives} (att);
\draw[arr] (att) -- (qfp);
\draw[arr] (qfp) -- (srch);
\draw[arr] (idx) -- (srch);
\node[lab, anchor=north] at (5.775,-2.55) {queries};

% lane 3: what is measured
\node[metric] (m1) at (0,-3.75) {\textbf{Robustness}\\\scriptsize R@1, $\mu$AP, TPR at calibrated FPR};
\node[metric] (m2) at (3.85,-3.75) {\textbf{Localization}\\\scriptsize edited region; time offset};
\node[metric] (m3) at (7.7,-3.75) {\textbf{Security}\\\scriptsize evasion, forgery, inversion};
\node[metric] (m4) at (11.55,-3.75) {\textbf{Complementarity}\\\scriptsize watermark key vs.\ fingerprint};
\draw[densely dashed, black!50, line width=0.5pt] ($(m1.north west)+(-0.1,0.18)$) -- ($(m4.north east)+(0.1,0.18)$);
\end{tikzpicture}

\caption{Benchmark protocol. \val{fp:count/methods} fingerprinting methods are applied to registered assets and distractors (registry) and to queries derived from registered assets (positives) or from never-registered media (negatives) under \val{fp:count/attacks} transformations and six kinds of partial edit. Thresholds are fixed on half of the negatives and evaluated on the other half. The same queries also drive the localization, security and watermark-complementarity measurements.}
\label{fig:fp-protocol}
\end{figure*}

\subsubsection{Methods Under Test}
We included every fingerprinting method we could find that is distributed with an implementation and, where learned, with weights, and that targets natural images, music and speech, or natural video (Table~\ref{tab:fp-methods}). We recorded the licence of the code and, separately, of the weights and of the data they were trained on, because the three differ in several projects. A method is \emph{product-eligible} (tier~P) when code and weights carry a permissive licence (MIT, BSD, Apache-2.0 or CC0) and the weights were not trained on data restricted to non-commercial use. The strongest learned copy detectors are not: SSCD and the ISC21 winner are MIT-licensed code, but their weights carry no licence of their own and were trained on DISC21, whose repository is CC BY-NC 4.0; DINOHash has no licence file; and the neural music fingerprint NMFP is GPL-3.0. We measured them as reference points (tier~R, marked \reftier), as the watermarking track does for its unlicensed checkpoints (Section~\ref{sec:wm-method}). We excluded Apple NeuralHash (weights extracted from a device, no licence), PhotoDNA (licensed only for child-safety use), Adobe's image comparator network (no public code), SuperPoint and TruFor (non-commercial weights), and Olaf and Panako, which we did not run: they are AGPL-3.0, whose network-use clause rules them out for a hosted provenance service, and Chromaprint and audfprint already represent their landmark and sub-band designs.

\paragraph{Images.} Twenty-one descriptors: PDQ \cite{meta2019pdq}, with and without matching the query's eight dihedral variants; aHash, dHash, pHash (64 and 256 bits) and wHash from imagehash \cite{buchner_imagehash}; BlockMean, Marr--Hildreth and ColorMoment from OpenCV \cite{opencv_imghash}; Blockhash \cite{commonsmachinery_blockhash}; the ISCC Image-Code at 64 bits (the ISO default) and 256 bits \cite{iso24138}; DINOv2-S and DINOv2-B class tokens \cite{oquab2024dinov2}; a keyed 256-bit sign projection of DINOv2-S (below); OpenCLIP ViT-B/32 \cite{cherti2023openclip}; SSCD with ResNet-50 and ResNeXt-101 backbones \cite{pizzi2022sscd}; the ISC21 descriptor-track winner \cite{yokoo2021isc}; and DINOHash \cite{singhi2025dinohash}. ISCC's experimental semantic image code \cite{iscc_sci} ran in our pilot but not at full scale: it binarises an ISC21-derived descriptor that we evaluate directly, and its reference model runs one image at a time at about 3\,s per image on a CPU core. Learned descriptors see the whole image resized to their square input without cropping, so that a crop attack cannot remove content the descriptor was never shown. Blockhash's reference implementation loops over pixels in Python; we vectorised it and verified bit-identical output on five image sizes including odd ones.

\paragraph{Keyed projection.} DINOv2-S LSH-256 multiplies the L2-normalised DINOv2-S descriptor by a $384 \times 256$ Gaussian matrix drawn from a secret seed and keeps the signs. It is a standard random-hyperplane hash, included because it is compact, product-eligible and, unlike a raw descriptor, meaningless to anyone without the key; the security experiments test that claim.

\paragraph{Audio.} Chromaprint \cite{lalinsky_chromaprint}, with candidate generation on the 20 most significant bits of each sub-fingerprint and verification by bit error rate at the voted offset (the AcoustID index design); the ISCC Audio-Code at 64 and 256 bits, computed from the whole-file Chromaprint vector; audfprint's spectral-peak landmarks \cite{ellis_audfprint}, scored by the number of hashes that agree on one time offset; LAION-CLAP \cite{wu2023clap} as a semantic embedding, mean-pooled over 10\,s windows; and NMFP \cite{araz2025nmfp}, a neural music fingerprint with 1\,s segment embeddings, run in its authors' TensorFlow environment and matched by segment voting.

\paragraph{Video.} vPDQ \cite{meta_vpdq} at one hash per second, scoring a candidate by the share of query frames with a reference frame within Hamming distance 31; TMK+PDQF \cite{meta2019pdq, poullot2015tmk}, scored by its level-2 time-kernel similarity using Meta's command-line tools; the ISCC Video-Code at 64 and 256 bits from MPEG-7 frame signatures \cite{paschalakis2012mpeg7}; and DINOv2-S and SSCD on frames sampled at 1\,fps, either mean-pooled into one clip descriptor or matched frame by frame with temporal voting, the baseline recipe of the 2023 Video Similarity Challenge \cite{pizzi2024vsc}.

\begin{table*}[!t]
\centering
\scriptsize
\caption{Methods, licences, stored size and extraction time. Bits: stored bits per registered asset (float descriptors count 32 bits per dimension; time-based systems report the mean over the registered items). Time: per image (1024 px), per second of audio or per 5\,s clip.}
\label{tab:fp-methods}
\setlength{\tabcolsep}{3pt}
\begin{tabular}{lllp{3.3cm}p{3.6cm}rr}
\toprule
Modality & Method & Class & Code licence & Weights / data & Bits & Time (ms) \\
\midrule
Image & PDQ & Perceptual hash & BSD (ThreatExchange); pdqhash binding  & n/a & 256 & 31 CPU \\
 & PDQ (dihedral) & Perceptual hash & BSD (ThreatExchange); pdqhash binding  & n/a & 256 & 31 CPU \\
 & aHash & Perceptual hash & BSD-2-Clause (imagehash) & n/a & 64.0 & 4.3 CPU \\
 & dHash & Perceptual hash & BSD-2-Clause (imagehash) & n/a & 64.0 & 4.4 CPU \\
 & pHash & Perceptual hash & BSD-2-Clause (imagehash) & n/a & 64.0 & 5.2 CPU \\
 & pHash-256 & Perceptual hash & BSD-2-Clause (imagehash) & n/a & 256 & 5.6 CPU \\
 & wHash & Perceptual hash & BSD-2-Clause (imagehash) & n/a & 64.0 & 63 CPU \\
 & BlockMean & Perceptual hash & Apache-2.0 (opencv\_contrib) & n/a & 256 & 1.0 CPU \\
 & Marr–Hildreth & Perceptual hash & Apache-2.0 (opencv\_contrib) & n/a & 576 & 5.5 CPU \\
 & ColorMoment & Perceptual hash & Apache-2.0 (opencv\_contrib) & n/a & 1344 & 3.6 CPU \\
 & Blockhash & Perceptual hash & MIT & n/a & 256 & 19 CPU \\
 & ISCC Image-64 & ISCC (ISO 24138) & Apache-2.0 & n/a & 64.0 & 11 CPU \\
 & ISCC Image-256 & ISCC (ISO 24138) & Apache-2.0 & n/a & 256 & 11 CPU \\
 & DINOv2-S & General embedding & Apache-2.0 & Apache-2.0 (HF model card) & 12288 & 2.0 GPU \\
 & DINOv2-B & General embedding & Apache-2.0 & Apache-2.0 (HF model card) & 24576 & 3.9 GPU \\
 & DINOv2-S LSH-256 & General embedding & Apache-2.0 & Apache-2.0 (HF model card) & 256 & 1.7 GPU \\
 & OpenCLIP B/32 & General embedding & MIT (open\_clip) & MIT (laion/CLIP-ViT-B-32-laion2B-s34B-b79K) & 16384 & 1.1 GPU \\
 & SSCD R50\reftier & Copy-detection trained & MIT & none stated; trained on DISC21 (CC BY-NC 4.0 & 16384 & 2.7 GPU \\
 & SSCD RX101\reftier & Copy-detection trained & MIT & none stated; trained on DISC21 (CC BY-NC 4.0 & 32768 & 4.6 GPU \\
 & ISC21 1st\reftier & Copy-detection trained & MIT & MIT repo release; trained on DISC21 (CC BY-N & 8192 & 8.3 GPU \\
 & DINOHash-96\reftier & Copy-detection trained & none (no LICENSE file) & none (no LICENSE file) & 96.0 & 3.6 GPU \\
\addlinespace
Audio & Chromaprint & Perceptual hash & MIT (chromaprint core; fpcalc binary l & n/a & 7072 & 2.0 CPU \\
 & audfprint & Perceptual hash & MIT & n/a & 40397 & 1.4 CPU \\
 & ISCC Audio-64 & ISCC (ISO 24138) & Apache-2.0 & n/a & 64.0 & 2.0 CPU \\
 & ISCC Audio-256 & ISCC (ISO 24138) & Apache-2.0 & n/a & 256 & 2.0 CPU \\
 & CLAP & General embedding & CC0-1.0 & CC0-1.0 (LAION-AI/CLAP release) & 16384 & 1.1 GPU \\
 & NMFP\reftier & Copy-detection trained & GPL-3.0 & Zenodo 15719945 (licence as repo, GPL-3.0);  & 239108 & -- \\
\addlinespace
Video & vPDQ & Perceptual hash & BSD (ThreatExchange vpdq) & n/a & 1258 & 132 CPU \\
 & TMK+PDQF & Perceptual hash & BSD (ThreatExchange tmk) & n/a & 2106752 & 51 CPU \\
 & ISCC Video-64 & ISCC (ISO 24138) & Apache-2.0 & n/a & 64.0 & 54 CPU \\
 & ISCC Video-256 & ISCC (ISO 24138) & Apache-2.0 & n/a & 256 & 54 CPU \\
 & DINOv2-S frames & General embedding & Apache-2.0 & Apache-2.0 (HF model card) & 61440 & 136 GPU \\
 & DINOv2-S mean & General embedding & Apache-2.0 & Apache-2.0 (HF model card) & 12288 & 136 GPU \\
 & SSCD frames\reftier & Copy-detection trained & MIT & none stated; trained on DISC21 (CC BY-NC 4.0 & 81920 & 137 GPU \\
 & SSCD mean\reftier & Copy-detection trained & MIT & none stated; trained on DISC21 (CC BY-NC 4.0 & 16384 & 137 GPU \\
\bottomrule
\end{tabular}
\par\vspace{2pt}\parbox{\linewidth}{\raggedright\scriptsize \reftier~Reference tier: weights without a licence, trained on non-commercial data, or copyleft code; measured for comparison, not product-eligible.}
\end{table*}


\subsubsection{Corpora}
\emph{Images.} Amazon Berkeley Objects (ABO, CC BY 4.0) \cite{collins2022abo} supplies product photographs from real marketplace listings, the setting in which copied and relabelled product images circulate. We sampled \val{fp:corpus/abo/reg} registered originals (one main image per product, shorter side at least 800\,px), stored at a longer side of 1024\,px as the asset a provenance service would register, and \val{fp:corpus/abo/dist} distractors (one image per other product). Queries come from \val{fp:corpus/abo/pos} registered originals; negative queries from \val{fp:corpus/abo/neg} images of products that are neither registered nor distractors; and hard negatives from \val{fp:corpus/abo/hard} \emph{other} photographs of registered products. ABO reuses photographs across listings and colour variants, so before scoring we flagged pairs that are copies of each other under a rule fixed in advance on unattacked originals (PDQ distance at most 16 bits, SSCD cosine at least 0.80, or DINOv2-B cosine at least 0.95). The rule removed \val{fp:dedup/n_dist_flagged} distractors and \val{fp:dedup/n_neg_flagged} negative sources, and marked \val{fp:dedup/n_hard_near_duplicate} of \val{fp:dedup/n_hard} hard negatives as recoloured versions of the registered photograph, which we report separately. As a second, literature-comparable set we fetched a subset of the DISC21 development split \cite{douze2021isc}: \val{fp:corpus/disc/gt} ground-truth query--reference pairs, \val{fp:corpus/disc/negq} queries without a reference, and \val{fp:corpus/disc/refs} references in total. DISC21 is licensed for non-commercial research; we use it for measurement only.

\emph{Audio.} \val{fp:corpus/audio/music} music tracks from the FMA large set \cite{defferrard2017fma}, restricted to tracks released under CC BY, CC BY-SA or CC0 and excluding the FMA-medium subset, and \val{fp:corpus/audio/speech} LibriSpeech utterances of at least 8\,s \cite{panayotov2015librispeech}, all at 22.05\,kHz mono. Music clips are 30\,s. Of these, \val{fp:corpus/audio/reg} are registered, \val{fp:corpus/audio/dist} are distractors and \val{fp:corpus/audio/neg} are negative sources.

\emph{Video.} Five-second clips cut every 5\,s from four Blender open films \cite{blender2010sintel, blender2008bbb, blender2012tos, blender2006ed}, skipping titles and end credits, and the first 5\,s of \val{fp:corpus/video/xiph} Xiph test sequences \cite{xiphderf}, scaled to 1280\,px wide at 25\,fps and stored as H.264 at CRF~12: \val{fp:corpus/video/reg} registered, \val{fp:corpus/video/dist} distractor and \val{fp:corpus/video/neg} negative clips. Clips from the same film share style and characters, which makes the distractors deliberately hard.

\subsubsection{Transformations}
\emph{Images} (\val{fp:count/image/attacks}, grouped into families in Fig.~\ref{fig:fp-image-heatmap}): JPEG at quality 75, 50 and 30 and WebP at 50; downscaling to 0.5 and 0.25 and a 0.75 aspect change; centred crops keeping 90, 70, 50 and 30\% of each side and an off-centre 50\% crop; rotation by 5, 15 and 90 degrees, horizontal flip, perspective, skew and 25\% padding; brightness, contrast, saturation, gamma, greyscale, blur, Gaussian noise and pixelation; text, emoji (30\% of the image), meme caption, screenshot and overlay onto another photograph at half size, following AugLy \cite{papakipos2022augly}; and two chains, a social-media chain (downscale, JPEG~70, 90\% crop) and a re-capture simulation (perspective, defocus, sensor noise, tone curve, JPEG~80). \emph{Audio} (\val{fp:count/audio/attacks}): MP3 and AAC at 64\,kb/s, Opus at 24\,kb/s, resampling through 8\,kHz; white noise at 20 and 10\,dB SNR and babble from other music at 10 and 0\,dB; telephone band, 2\,kHz low-pass, dynamic compression and reverberation; pitch $+2$ semitones, tempo 0.9 and 1.1 and speed 1.05; excerpts of 10 and 5\,s at a random offset, an excerpt with MP3 and noise, and a re-recording simulation. Unlike the watermarking track, codec outputs are not re-aligned: a fingerprint must cope with codec delay itself. \emph{Video} (\val{fp:count/video/attacks}): H.264 at CRF~28 and 35, H.265 at CRF~32; downscaling to 0.5 and 0.25; 75\% and 50\% crops; flip, 90 and 5 degree rotation, letterboxing; greyscale, brightness, contrast, blur and noise; text caption, logo, and picture-in-picture at half size on another clip; 15\,fps, 1.25$\times$ speed, a 2.5\,s excerpt, the same 2.5\,s inserted into 5\,s of another clip, and a screen-recording simulation.

\subsubsection{Retrieval Metrics and Operating Points}
Search is exact (flat index, top~50): Hamming similarity for binary codes, cosine for L2-normalised descriptors, and each audio or video system's own score. We report Recall@1 and the micro-average precision ($\mu$AP) of the ISC21 protocol \cite{douze2021isc}, both threshold-free. A registry, however, must decide, and a false decision binds content to someone else's manifest. We therefore fix thresholds on the \emph{calibration} half of the negative sources (all their transformed copies) and report on the other half. Two kinds of target are used. The query-level target fixes the rate at which a never-registered query's best match passes at the benchmark's registry size (1\%). The pair-level target fixes the rate at which a (negative query, reference) pair passes. For a registry of $N$ assets and an average pair-level false-match probability $p$, the expected number of false matches per negative query is $Np$; the probability of at least one false match is at most $\min(1,Np)$, without requiring independent reference comparisons. Our headline target is a pair-level rate of $10^{-7}$: at the benchmark's \val{fp:count/image/nref} references, this corresponds to about one expected false match per hundred negative queries in aggregate. This calculation is not a measured query-level false-binding rate. Table~\ref{tab:fp-ci} reports held-out query-level rates at the separately calibrated query-level threshold. We also report $10^{-8}$, the smallest rate our \val{fp:count/image/neg_cal} calibration queries resolve. A positive query counts as detected only if its best match is the true asset and passes the threshold. A target is reported only if the calibration set resolves it, that is, if it implies at least five false pairs (or queries) among the calibration negatives; otherwise the threshold would merely sit above every negative observed. This is a reporting-resolution rule, not evidence that the comparisons are independent or that the tail probability is estimated precisely. The queries are not independent: every registered or negative source contributes one query per transformation, and the transformed copies of one source tend to succeed or fail together. Intervals for these retrieval metrics therefore come from a bootstrap over sources (2{,}000 replicates): registered sources are drawn with replacement together with all their queries, calibration negative sources likewise, with the threshold recomputed from the drawn calibration set in every replicate, and held-out negative sources for the out-of-sample false-match rate. Calibration and held-out negatives come from disjoint sources. Planned paired comparisons use a sign-flip test on each source's difference in detections between the two methods, with Holm adjustment \cite{holm1979}, and report a bootstrap interval of the difference in detection rate; both methods are resampled with the same sources. The sign-flip test and the verification experiment (Table~\ref{tab:fp-verify}) hold the thresholds at their full-sample values, so their uncertainty excludes threshold estimation; the bootstrap interval of a paired difference recomputes the thresholds in every replicate.

\subsubsection{Partial Edits and Localization}
Partial edits are generated from \val{fp:edit/n_sources} registered originals, with the edited region known exactly: a smooth blob of 1, 3, 10 or 25\% of the image, placed where the image has texture. The edit is a splice from another photograph, a copy-move within the image, classical (Telea) or learned (LaMa \cite{suvorov2022lama}) inpainting of the region, a hue rotation, or replacement by a flat label with new text, followed by no post-processing, JPEG~75, downscaling with JPEG~80, or a 90\% crop. A balanced design gives \val{fp:edit/n_edits} edited images. We first ask whether each fingerprint still retrieves the original at its calibrated threshold. We then localize the edit against the registered original: the query is registered onto the original with SIFT and RANSAC \cite{lowe2004sift, fischler1981ransac} or with DISK and LightGlue \cite{tyszkiewicz2020disk, lindenberger2023lightglue}, and compared through six difference maps: blurred pixel difference, SSIM \cite{wang2004ssim}, LPIPS \cite{zhang2018lpips}, DINOv2-S patch distance, PDQ over an $8 \times 8$ grid of tiles, and a fusion of the five. Map thresholds are fixed at a pixel false-positive rate of 1\% on unedited originals put through the same post-processing, on half of the sources; the other half is evaluated. For audio and video we measure temporal localization: how accurately each system recovers the offset of an excerpt within its source, and which interval of an inserted clip matches.

\subsubsection{Security}
We consider an adversary who knows the fingerprinting method and wants either to \emph{evade} it (republish a registered asset so that it no longer matches) or to \emph{forge} a match (make an unrelated image bind to a registered asset, for example to borrow its credentials). Evasion uses projected gradient descent \cite{madry2018pgd} under an $L_\infty$ budget of 2, 4, 8 or 16 grey levels: white-box on the learned descriptors, through differentiable re-implementations of PDQ, pHash, aHash, dHash and the ISCC Image-Code for those hashes, and by transfer from an ensemble of DINOv2-S, SSCD and OpenCLIP for all methods. Success is always judged by the real extractor on the 8-bit result. Image and audio attacks were judged at the method's pair-level $10^{-6}$ threshold, and video attacks at its query-level 1\% threshold, because the video calibration set does not resolve $10^{-6}$ (Section~\ref{sec:fp-res-video}). For images, $10^{-6}$ is a looser operating point than the headline: an evasion must push the score below it, which makes the reported evasion rates conservative, and we re-score white-box evasions at the $10^{-7}$ threshold from their final scores. Forgery success at $10^{-6}$ is correspondingly an upper bound on success at $10^{-7}$. Forgery pushes a negative image towards a registered target and succeeds only if the target becomes the top match in the full registry above threshold. Because a soft-binding value is written into the manifest, we also test \emph{inversion}: a decoder trained on distractor descriptors reconstructs a 64$\times$64 thumbnail from a registered asset's descriptor, and we measure whether the reconstruction identifies the asset among the others. For audio and video we apply the same attacks to the learned systems and measure, for every method, the smallest pitch, tempo, speed or crop change that breaks it.

\subsubsection{Watermark Complementarity}
The registered assets of the three modalities were also watermarked with MIT-licensed models from the watermarking track: TrustMark~Q at its default strength for images \cite{bui2025trustmark}, AudioSeal \cite{sanroman2024audioseal} and Video Seal \cite{fernandez2024videoseal}. The marked copies went through the same transformations. For every query we record whether the watermark satisfies its method-specific success criterion, and whether each fingerprint identifies the asset. The three watermarks decide this differently, and the decisions relate differently to recovering a registry identifier. TrustMark transmits a 100-bit codeword; in the BCH\_5 mode we use it carries a 61-bit data payload, 35 parity bits and 4 version bits. Its BCH decoder either corrects the codeword or reports failure; the key counts as recovered when the decoded 61-bit payload equals the asset's key exactly, which is identification: the decoded payload itself names the registry entry. AudioSeal (the base model) carries only 16 bits and has no error correction. Its detection score is the fraction of audio frames classified as watermarked; we require at least half of the frames, a more permissive rule than the 0.8 of the watermarking track's blind detector (Section~\ref{sec:wm-fp-method}), because here detection counts only together with an exact match of all 16 bits to the asset's key. No AudioSeal decode of the watermarking track's unmarked audio exceeded \val{wm:audioseal/unmarked_max_score}, so its blind false-positive count is the same under either rule. A recovered key identifies an asset only if no other asset holds it. The space contains only $2^{16}$ keys, so a registry larger than that cannot give every asset a distinct AudioSeal key; our keys were drawn pseudo-randomly per asset without a uniqueness check, and the \val{fp:wm/audio/n_assets} registered assets received \val{fp:wm/audio/nkeys} distinct keys, so \val{fp:wm/audio/n_colliding} assets share one key. Video Seal carries 256 bits without error correction, and exact recovery of all 256 bits is rare after compression; we count the query as passing expected-key verification when at least $k$ of the 256 decoded bits agree with the asset's key, with $k$ the smallest value for which a random 256-bit word agrees with a given key with probability at most $10^{-6}$. For positives this is verification against the expected key, not a test of unique identification among all registered keys. A registry reader without that expectation would compare the decoded bits with every registered key; on the never-registered negatives we did exactly that, so a false Video Seal key there means agreement with \emph{any} of the registered keys. Under the independent-fair-bit null used to set $k$, the union bound gives a per-query false-match probability of at most $\min(1,N \times 10^{-6})$ for $N$ keys. This analytic bound is conditional on that null and does not replace the measured false-match rate. Testing all keys on negative queries does not establish unique identification on positive queries. For AudioSeal negatives we report both the detector's false detections and detections whose message equals a registered key. We also measure how far the watermark itself moves each fingerprint, and whether the registry should hold the fingerprint of the original or of the marked asset.

\subsubsection{Second-Stage Verification}
To test whether checking a retrieved candidate against the stored original removes the false matches that a global descriptor lets through, we ran SIFT with RANSAC between every image query (positive and negative) and the registered original of its top-1 candidate under DINOv2-S and PDQ, at the localization track's working size, and recorded the number of geometric inliers. A match is accepted when the retrieval score passes a threshold and the inlier count reaches $k$; both are fixed on the calibration negatives and evaluated on the held-out ones.

\subsubsection{Hosts and Reproducibility}
Learned descriptors were extracted on one NVIDIA A10G (AWS g5.2xlarge, AMD EPYC 7R32); CPU descriptors single-threaded on the same processor model (an AWS c5a.16xlarge was added to spread the CPU work). Every source tree, checkpoint and corpus file is pinned by commit or SHA-256. Because the watermarking track found host-dependent decoder outputs (Section~\ref{sec:wm-res-stability}), we fingerprinted the same 200 PNG files on the GPU host and on an Apple M5 Pro laptop and compared the descriptors bit by bit.

The retrieval tracks and the watermark combination were run twice, on 2026-09-26 and again on 2026-09-28 on fresh hosts of the same types, because the first run did not keep the per-query scores that source-level intervals require. The corpus was rebuilt from its sources and matched the pinned SHA-256 of every image, audio and video item. The image track reproduced the first run exactly (every pooled rate to within $10^{-6}$), and the video track to within two of 3{,}000 queries. The audio track did not reproduce exactly, for two reasons that we found and fixed: the background music of the two babble transformations had been drawn from whichever surplus FMA downloads were present, which differed between runs (the pool is now pinned by name and hash), and CLAP was extracted on the CPU host in the second run instead of the GPU. All reported numbers come from the second run.

